\chapter{Introduction}
\label{ch:intro}

This is an example chapter. Abbreviations are typeset with \gls{ai} and
\gls{ml}; symbols such as \gls{sym:E} come from \texttt{symbols.tex}.
References are cited as usual~\cite{lamport1994}.

Figures, tables and algorithms are collected by \texttt{endfloat} and
printed in the appendices at the end of the document
(see Figure~\ref{fig:example}, Table~\ref{tab:example} and
Algorithm~\ref{alg:example}).

\begin{figure}[htbp]
  \centering
  \includegraphics[width=0.5\textwidth]{example-image}
  \caption{An example figure.}
  \label{fig:example}
\end{figure}

\begin{table}[htbp]
  \centering
  \caption{An example table.}
  \label{tab:example}
  \begin{tabular}{lr}
    \toprule
    Item & Value \\
    \midrule
    A & 1 \\
    B & 2 \\
    \bottomrule
  \end{tabular}
\end{table}

\begin{algorithm}[htbp]
  \caption{An example algorithm.}
  \label{alg:example}
  \begin{algorithmic}[1]
    \For{$i \gets 1$ to $n$}
      \State $s \gets s + i$
    \EndFor
  \end{algorithmic}
\end{algorithm}

\begin{definition}
  This is an example definition environment.
\end{definition}

\lipsum[1-3]
